\section{Fractal observables from the largest stabilized network}
\label{sec:fractal-counts-methods}

The fractal observables were measured on stabilized samples of the
self-organized percolation process.  For each independent seed, after the
stabilization time had elapsed, the simulation collected \(M\) independent
windows separated by a prescribed gap.  Each collected window defines one
network sample of linear size \(L\).  The standard observables already present
in the simulation pipeline were kept unchanged.  In addition, for the largest
systems, namely \(L=16384\) in two dimensions and \(L=1024\) in three
dimensions, geometric counting observables were computed for each collected
network.

All measurements described below were performed on the largest connected
component of each stabilized network.  For node percolation, connectivity is
defined by nearest-neighbor occupied sites.  For bond percolation, two occupied
sites are connected only when the bond between them is active.  The transverse
directions are periodic, while the growth direction is open.  In two
dimensions this corresponds to periodic boundary conditions along \(x\) and
open boundary conditions along \(y\).  In three dimensions the \(x\) and \(y\)
directions are periodic and the \(z\) direction is open.

\subsection{Mass fractal dimension}
\label{subsec:mass-fractal-dimension}

The mass fractal dimension \(d_f\) was estimated by a box-counting procedure
applied to the sites of the largest connected component.  For a box size
\(\epsilon\), the number of boxes intersecting the largest component is denoted
by \(N_{\mathrm{bulk}}(\epsilon)\).  Several offsets of the box grid were used
for each value of \(\epsilon\), and the stored count corresponds to the average
over offsets for that sample.

The expected scaling form is
\begin{equation}
    N_{\mathrm{bulk}}(\epsilon)
    \sim
    \left(\frac{L}{\epsilon}\right)^{d_f}.
    \label{eq:df-box-counting}
\end{equation}
Thus, \(d_f\) is obtained as the slope of a linear regression of
\(\log N_{\mathrm{bulk}}\) versus \(\log(L/\epsilon)\).  The processed data
store, for every network sample, the pair of arrays
\(\{L/\epsilon, N_{\mathrm{bulk}}(\epsilon)\}\), so that the fit range can be
chosen at the analysis stage.

\subsection{Hull and external hull dimensions}
\label{subsec:hull-dimensions}

The complete hull and the external hull were also analyzed by box counting.
The complete hull contains all boundary elements of the largest connected
component, including boundaries surrounding internal voids.  The external hull
contains only the boundary elements that are accessible from the exterior of
the component.  Therefore, internal cavities do not contribute to the external
hull.

For a box size \(\epsilon\), let \(N_{\mathrm{hull}}(\epsilon)\) be the number
of boxes intersecting the complete hull, and let
\(N_{\mathrm{hull,ext}}(\epsilon)\) be the corresponding number for the
external hull.  The scaling relations are
\begin{align}
    N_{\mathrm{hull}}(\epsilon)
    &\sim
    \left(\frac{L}{\epsilon}\right)^{d_{\mathrm{hull}}},
    \label{eq:dhull-box-counting}
    \\
    N_{\mathrm{hull,ext}}(\epsilon)
    &\sim
    \left(\frac{L}{\epsilon}\right)^{d_{\mathrm{hull,ext}}}.
    \label{eq:dhullext-box-counting}
\end{align}
As in the mass measurement, the simulation stores the offset-averaged
box-counting curves for each network sample.  The exponents
\(d_{\mathrm{hull}}\) and \(d_{\mathrm{hull,ext}}\) are obtained as the slopes
of the respective log-log plots.

\subsection{Chemical distance exponent}
\label{subsec:chemical-distance-exponent}

The chemical distance exponent \(d_{\min}\) was measured from shortest paths
inside the largest connected component.  For each sample, a set of origin sites
was randomly selected from the largest component.  A breadth-first search was
then performed from each origin using the same connectivity rule as the
percolation model.  For every reached site, the algorithm records two
distances:
\begin{itemize}
    \item the Euclidean distance \(r\) from the origin, measured with periodic
    wrapping in the transverse directions only;
    \item the chemical distance \(\ell\), defined as the graph distance from
    the origin along the connected component.
\end{itemize}

Distances were grouped into logarithmic bins in \(r\).  For each bin, the
processed data store the minimum chemical distance
\(\ell_{\min}(r)\), the mean chemical distance \(\ell_{\mathrm{mean}}(r)\),
the maximum chemical distance \(\ell_{\max}(r)\), the standard deviation, and
the number of contributing pairs.  The default estimator used for
\(d_{\min}\) is \(\ell_{\min}(r)\), using all available bins.  The mean and
maximum values are kept only as diagnostic quantities.

The scaling relation used in the analysis is
\begin{equation}
    \ell_{\min}(r) \sim r^{d_{\min}}.
    \label{eq:dmin-chemical-distance}
\end{equation}
Consequently, \(d_{\min}\) is obtained as the slope of a linear regression of
\(\log \ell_{\min}\) versus \(\log r\).  No rescaling by \(L\) is applied to
this observable.  The quantity \(L/\epsilon\) is used only for box-counting
observables.

The breadth-first search may optionally be limited by a maximum chemical
distance \(\ell_{\max}^{\mathrm{BFS}}\).  The processed files therefore also
store a boolean flag indicating whether a radial bin reached this cutoff.  This
flag is not used by default in the fit, but it is useful for diagnosing
finite-search effects or for testing alternative fitting windows.

\subsection{Minimum-path yardstick estimator}
\label{subsec:minimum-path-yardstick}

In addition to the radial chemical-distance estimator, the code also stores a
yardstick estimator applied directly to one shortest path connecting the bottom
and top boundaries.  For each post-stabilization sample, a multi-source
breadth-first search is initialized from all sites of the largest component
touching the bottom boundary.  The search uses the same node or bond
connectivity rule as the simulation and stops when the first site on the top
boundary is reached.  This gives one shortest bottom-to-top path with chemical
length \(\ell_{\mathrm{bt}}\).

The full path is not stored.  Instead, for each yardstick radius \(R\), the
algorithm walks along the reconstructed path.  Starting from the current center,
it advances along the path until the Euclidean distance from the center is at
least \(R\), counts one yardstick sphere, and uses that path site as the next
center.  A final partial segment is counted once when the end of the path is
reached before another full radius \(R\) can be placed.  The stored curve is
\begin{equation}
    N(R) \sim \left(\frac{L}{R}\right)^{d_{\min}^{\mathrm{yardstick}}}.
    \label{eq:dmin-yardstick}
\end{equation}
Thus the yardstick estimate is obtained as the slope of
\(\log N(R)\) versus \(\log(L/R)\).  The output stores only \(R\), \(L/R\),
\(N(R)\), the path length, and the two path endpoints.

\subsection{Processed data format}
\label{subsec:processed-counts-format}

The raw count files are stored separately from the standard processed
observables.  The processed fractal-count data are organized by percolation
type, dimension, system size, threshold parameter, control parameter, and
density:
\begin{verbatim}
SOP_data/published_counts/
  <type>_percolation/
    dim_<dim>/
      L_<L>/
        fT_<f_T>/
          c_<c>/
            rho_<rho>/
              counts_P0_<P0>_p0_<p0>.json
\end{verbatim}
The values of \(P_0\) and \(p_0\) appear in the processed filename.  Each JSON
file contains a \texttt{meta} section with the simulation parameters and a
\texttt{samples} list.  Each entry of \texttt{samples} corresponds to one
collected network and contains the arrays required for the log-log regressions:
\begin{align*}
    d_f &: \quad
    \{L/\epsilon,\; N_{\mathrm{bulk}}(\epsilon)\},
    \\
    d_{\mathrm{hull}} &: \quad
    \{L/\epsilon,\; N_{\mathrm{hull}}(\epsilon)\},
    \\
    d_{\mathrm{hull,ext}} &: \quad
    \{L/\epsilon,\; N_{\mathrm{hull,ext}}(\epsilon)\},
    \\
    d_{\min} &: \quad
    \{r,\; \ell_{\min}(r)\}.
\end{align*}
The exponent estimates are not fixed during data processing.  Instead, the
processed files preserve the scale-dependent curves for each network sample,
allowing the fitting interval and estimator to be chosen later in the notebook
analysis.
